% ============================================================
%  ABSTRACT
% ============================================================
\begin{abstract}

\ac{FAST} is a rapid bedside ultrasound protocol that is critical in
emergency medicine for detecting life-threatening internal injuries.
Performing \ac{FAST} scans autonomously with a mobile robot could extend
emergency care to mass-casualty events, disaster sites, and scenarios where
trained sonographers are unavailable. However, autonomous \ac{FAST} scanning
on a mobile platform poses two fundamental challenges: the robot must navigate
to an unknown patient location using uncertain visual detection, and must
simultaneously position the ultrasound probe with millimetre-level accuracy
over the target anatomy.

This paper presents an integrated system for autonomous \ac{FAST} scanning
mounted on a Boston Dynamics Spot legged robot equipped with a Unitree Z1
manipulator arm. We formulate \ac{WBC} as a single \ac{QP} that
simultaneously drives the non-holonomic Spot base and the six-\ac{DOF} Z1
arm, avoiding the sequential stop-then-manipulate paradigm common in existing
mobile manipulation literature. The holistic Jacobian couples base and arm
velocities through a manipulability-weighted damped least-squares formulation,
automatically redistributing effort from the arm to the base near singular
configurations. Patient localisation relies on a dual-camera pipeline: an
Orbbec \ac{RGBD} camera provides wide-field \ac{YOLO}-based skeleton
detection with \ac{KF}-filtered 3D torso tracking during navigation, while
an Intel RealSense camera mounted at the \ac{EE} takes over for fine-grained
surface estimation during scanning. A six-state \ac{FSM} orchestrates the
mission phases---approach, sensor handoff, scanning, and reactive workspace
extension---ensuring safe and reliable operation across a variety of patient
positions and orientations.

Experiments on a physical Spot+Z1 platform demonstrate that the proposed
\ac{WBC} reduces mission completion time by \XX\,\% compared to a
stop-and-manipulate baseline, while maintaining a probe positioning error
below \XX\,mm. To the best of our knowledge, this is the first system to
demonstrate fully autonomous \ac{FAST} scanning on a legged mobile manipulator.

\end{abstract}
